\label{chap:83-21}

\input{ampliaciones_sucesoras_20260824/parte_i_ii/owners/04ac_geometria_proyectiva_bloch_rev2}

\section{Construcción del espacio de Hilbert a partir de la conexión nonádica y sus operadores de fase}

\subsection{Realización matricial de cada profundidad}

Para \(d\ge1\), definimos
\[
  \mathcal C_9^{(d)}=(\mathbb Z/9\mathbb Z)^d,
  \qquad
  H_d=\ell^2(\mathcal C_9^{(d)}),
  \qquad
  A_d=B(H_d)\cong M_{9^d}(\mathbb C).
\]
La base canónica \(\{\delta_x:x\in\mathcal C_9^{(d)}\}\) conserva la
etiqueta discreta de cada estado. La medida uniforme induce la traza
normalizada
\[
  \tau_d(a)=\frac{1}{9^d}\Tr(a).
\]

La factorización
\[
  \mathcal C_9^{(d+1)}
  =
  \mathcal C_9^{(d)}\times\mathbb Z/9\mathbb Z
\]
produce una identificación unitaria
\[
  H_{d+1}\cong H_d\otimes\mathbb C^9.
\]
Ésta es la entrada exacta de la torre operatorial. No se identifica por ello
\((\mathbb Z/9)^3\) con \(\mathbb F_3^6\) como grupos: comparten \(729\)
elementos, pero sus operaciones y cociclos de acarreo son distintos.

\subsection{Representación de Weyl–Heisenberg de la fase y el desplazamiento sobre la fibra trítica}

Sea \(\omega=e^{2\pi i/9}\). Sobre
\(H_1=\ell^2(\mathbb Z/9\mathbb Z)\), definimos
\[
  (Xf)(r)=f(r-1),
  \qquad
  (Zf)(r)=\omega^r f(r).
\]
Entonces
\[
  X^9=Z^9=I,
  \qquad
  ZX=\omega XZ.
\]
El grupo generado por \(X,Z,\omega I\) es el grupo de Heisenberg finito de
orden \(9^3=729\).

\begin{theorem}[Stone--von Neumann finito]
Toda representación unitaria irreducible del grupo de Heisenberg finito en
la que el centro actúa mediante el carácter primitivo
\(\omega I\) es unitariamente equivalente a la representación anterior.
\end{theorem}

\begin{proof}
Sea \(v\) un autovector de \(Z\). Los vectores
\[
  v,Xv,\ldots,X^8v
\]
tienen autovalores distintos de \(Z\), pues
\(ZX^kv=\omega^kX^kZv\). Son ortogonales y generan un subespacio
invariante. Por irreducibilidad forman una base de toda la representación.
En esa base, \(X\) es el desplazamiento cíclico y \(Z\) la fase diagonal.
\end{proof}

\begin{controlbox}
La unicidad falla si no se fija el carácter central primitivo o si se
permiten representaciones reducibles. Este teorema finito no debe confundirse
con la versión continua para las relaciones canónicas de conmutación.
\end{controlbox}

\subsection{Proyectores aritméticos e incompatibilidad}

Las hojas aditiva y multiplicativa de \(\APP\) definen subespacios de
información y sus proyectores ortogonales
\(P_{\Sigma,d},P_{\Pi,d}\). Las identidades construidas satisfacen
\[
 \|[P_{\Sigma,2},P_{\Pi,2}]\|
 =\frac{\sqrt2}{3},
 \qquad
 \|[P_{\Sigma,3},P_{\Pi,3}]\|
 =\frac{\sqrt3}{4}.
\]
La no conmutatividad no se deduce de una analogía cuántica: es un teorema
matricial finito. Su lectura «pre-Heisenberg» es posterior a ese cálculo.

\begin{proposition}
Los proyectores anteriores no pertenecen conjuntamente a una misma
subálgebra booleana de proyecciones.
\end{proposition}

\begin{proof}
En una subálgebra abeliana todos los proyectores conmutan. Los conmutadores
mostrados son no nulos.
\end{proof}

\subsection{Estructura hilbertiana de la publicación matricial y memoria genealógica preservada}

El espacio de Hilbert proporciona superposición lineal, proyección ortogonal
y cálculo espectral. HMT aporta el índice genealógico de la base y distingue
qué operador representa una transición, una medida o una sombra. La
realización no sustituye el \(\TPK\): lo representa.

\begin{sourcebox}
Los proyectores APP y el par Weyl--Heisenberg forman parte de la construcción
precedente. La organización \(H_d,A_d\) y la prueba de unicidad se establecen
aquí mediante la aplicación del teorema clásico correspondiente a los objetos
HMT.
\end{sourcebox}


\section{Las \(9\) fases del TPK y la geometría continua}

\subsection{Separación tipada entre la elevación UHF de los rangos nonádicos y la prolongación geométrica de la rama marcada}

Sea \(p_j=|e_j\rangle\langle e_j|\) el proyector sobre una puerta marcada.
La factorización \(H_{d+1}\cong H_d\otimes\mathbb C^9\) admite dos
elevaciones:
\[
  \iota_d(a)=a\otimes I_9,
  \qquad
  j_{d,j}(a)=a\otimes p_j.
\]

\begin{proposition}[Unidad y traza]
Para todo \(a\in A_d\),
\[
  \tau_{d+1}(\iota_d(a))=\tau_d(a),
  \qquad
  \iota_d(I)=I,
\]
mientras
\[
  \tau_{d+1}(j_{d,j}(a))=\frac1{9}\tau_d(a),
  \qquad
  j_{d,j}(I)\ne I.
\]
\end{proposition}

\begin{proof}
Como \(\Tr(I_9)=9\) y \(\Tr(p_j)=1\),
\[
 \tau_{d+1}(a\otimes I_9)
 =\frac{\Tr(a)\,9}{9^{d+1}}
 =\tau_d(a),
\]
y
\[
 \tau_{d+1}(a\otimes p_j)
 =\frac{\Tr(a)}{9^{d+1}}
 =\frac{\tau_d(a)}9.
\]
Las identidades de unidad son inmediatas.
\end{proof}

\begin{intuitionbox}
«Mirar las nueve fases del TPK» significa conservar el operador identidad sobre la
fibra completa. Una puerta aislada conserva una historia marcada; las nueve
juntas conservan además la unidad y la normalización global.
\end{intuitionbox}

\subsection{El límite UHF}

\begin{theorem}[Cierre UHF nonádico]
El límite inductivo unital
\[
  A_{9^\infty}
  :=
  \varinjlim(A_d,\iota_d)
  \cong
  \bigotimes_{n=1}^{\infty}M_9(\mathbb C)
\]
es el álgebra UHF de tipo sobrenatural \(9^\infty=3^\infty\). Posee una
traza producto \(\tau\) compatible con todas las \(\tau_d\).
\end{theorem}

\begin{proof}
Las inclusiones son inyectivas, unitales y traza-preservantes. La unión
algebraica está formada por matrices finitas anidadas; su cierre normado es,
por definición, una UHF. La traza producto es la única compatible con las
trazas matriciales de todos los factores \citep{glimm1960}.
\end{proof}

\begin{theorem}[Factor hiperinfinito]
Sea \((\pi_\tau,\mathcal H_\tau,\Omega_\tau)\) la representación GNS de la
traza. Entonces
\[
  \pi_\tau(A_{9^\infty})''\cong\Rhyper,
\]
el factor hiperinfinito de tipo \(II_1\).
\end{theorem}

\begin{proof}
El cierre es finito porque \(\tau(I)=1\). La extremalidad de la traza lo hace
factor. Es hiperfinito porque la unión creciente de las matrices
\(\pi_\tau(A_d)\) es débilmente densa. La unicidad del factor separable
hiperfinito \(II_1\) completa la identificación
\citep{murrayvonneumann1936,connes1976}.
\end{proof}

\subsection{La rama marcada cambia de tipo}

Sea \(V_{d,j}:H_d\to H_{d+1}\), \(V_{d,j}\xi=\xi\otimes e_j\).
Entonces
\[
  j_{d,j}(a)=V_{d,j}aV_{d,j}^*.
\]
En el límite de los espacios, la unión de las esquinas contiene los
operadores de rango finito. Su cierre normado es
\(\mathcal K(H_\infty)\), y el bicommutante natural es
\(B(H_\infty)\), de tipo \(I_\infty\).

\begin{figure}[H]
\centering
\resizebox{\textwidth}{!}{%
\begin{tikzpicture}[>=Latex,node distance=12mm and 17mm,
  box/.style={rounded corners,draw=hmtblue,fill=hmtlightblue,
    minimum width=34mm,minimum height=13mm,align=center},
  gold/.style={rounded corners,draw=hmtgold,fill=hmtlightgold,
    minimum width=34mm,minimum height=13mm,align=center}]
  \node[box] (a) {\(M_{9^d}\)\\\scriptsize estado finito};
  \node[box,above right=of a] (full) {\(a\otimes I_9\)\\\scriptsize nueve fases del TPK};
  \node[gold,below right=of a] (one) {\(a\otimes p_j\)\\\scriptsize puerta marcada};
  \node[box,right=of full] (uhf) {\(\mathrm{UHF}(9^\infty)\)\\\scriptsize unital/tracial};
  \node[box,right=of uhf] (ii) {\(\Rhyper\)\\\scriptsize tipo \(II_1\)};
  \node[gold,right=of one] (k) {\(\mathcal K(H_\infty)\)\\\scriptsize compactos};
  \node[gold,right=of k] (i) {\(B(H_\infty)\)\\\scriptsize tipo \(I_\infty\)};
  \draw[->,very thick,hmtblue] (a)--(full);
  \draw[->,very thick,hmtblue] (full)--(uhf);
  \draw[->,very thick,hmtblue] (uhf)--(ii);
  \draw[->,very thick,hmtgold] (a)--(one);
  \draw[->,very thick,hmtgold] (one)--(k);
  \draw[->,very thick,hmtgold] (k)--(i);
  \node[above=2mm of full,text=hmtgreen]{\scriptsize unidad y traza};
  \node[below=2mm of one,text=hmtred]{\scriptsize traza \(\div9\)};
\end{tikzpicture}
}
\caption{Una puerta marcada y las nueve fases del TPK producen cierres de tipos
distintos.}
\label{p12-83-c21-s2:fig:operator-gates}
\end{figure}

\subsection{Dimensión continua desde rangos nonádicos}

Para una proyección \(P\in M_{9^d}\),
\[
  d_d(P)=\tau_d(P)=\frac{\rank(P)}{9^d}.
\]
Los valores pertenecen a
\(\mathbb Z[1/9]\cap[0,1]\), conjunto denso en \([0,1]\).

\begin{theorem}[Geometría continua discreta]
Para cada \(t\in[0,1]\) existe una proyección \(P_t\in\Rhyper\) con
\(\tau(P_t)=t\). Además,
\[
  P\sim Q
  \quad\Longleftrightarrow\quad
  \tau(P)=\tau(Q),
\]
donde \(\sim\) es equivalencia de Murray--von Neumann.
\end{theorem}

\begin{proof}
Sea \(m_d=\lfloor9^dt\rfloor\). Como
\(0\le m_{d+1}-9m_d\le8\), pueden elegirse proyecciones anidadas de rango
\(m_d\). Su límite fuerte \(P_t\) satisface, por normalidad de la traza,
\[
  \tau(P_t)=\lim_d\frac{m_d}{9^d}=t.
\]
La clasificación de proyecciones en un factor finito da la equivalencia
final.
\end{proof}

La dimensión continua no es cardinalidad fraccionaria. Mide equivalencia de
subespacios en el cierre. Cada aproximante conserva rango entero; la
continuidad aparece al cerrar toda la genealogía unital.

\begin{resultbox}
La conservación holográfica de la unidad se convierte en un criterio
operatorial falsable: sustituir \(I_9\) por \(p_j\) divide la traza y cambia
el tipo del límite. El conjunto de las puertas no es un censo; es la
condición de la inclusión unital y tracial.
\end{resultbox}


\section{Sistema dinámico profinito: traslación, álgebra lógica y espacio de estados}

\subsection{La 2.ª ruta al factor hiperinfinito}

La construcción HMT produce el grupo compacto
\[
  K_\infty\cong C_4\times\mathbb Z_3
\]
como límite de kernels de norma de Pell, junto con una traslación
\(T_u(x)=x+u\) minimal y únicamente ergódica para la medida de Haar
\(\mu_H\). Consideramos
\[
  M_{\rm clk}
  =
  L^\infty(K_\infty,\mu_H)\rtimes_{T_u}\mathbb Z.
\]

\begin{theorem}[Factor del reloj]
Si \(T_u\) es libre, ergódica y preserva Haar, entonces
\[
  M_{\rm clk}\cong\Rhyper.
\]
\end{theorem}

\begin{proof}
La acción libre y ergódica p.m.p. produce un factor finito con traza
\[
 \tau_{\rm clk}\!\left(\sum_nf_nV^n\right)
 =\int f_0\,d\mu_H.
\]
La amenabilidad de \(\mathbb Z\) hace hiperfinita la relación de órbitas;
por tanto el factor es el hiperinfinito \(II_1\).
\end{proof}

\begin{figure}[H]
\centering
\resizebox{\textwidth}{!}{%
\begin{tikzpicture}[>=Latex,node distance=12mm and 15mm,
  box/.style={rounded corners,draw=hmtblue,fill=hmtlightblue,
    minimum width=36mm,minimum height=13mm,align=center},
  boundary/.style={rounded corners,draw=hmtred,dashed,fill=white,
    minimum width=36mm,minimum height=13mm,align=center}]
  \node[box] (tower) {\((\mathbb Z/9)^d\)\\\(M_{9^d}\)};
  \node[box,right=of tower] (uhf) {\(\mathrm{UHF}(9^\infty)\)};
  \node[box,right=of uhf] (r1) {\(\Rhyper\)};
  \node[box,below=18mm of tower] (clock) {\(C_4\times\mathbb Z_3\)\\\(T_u\)};
  \node[box,right=of clock] (cross) {\(L^\infty(K_\infty)\rtimes\mathbb Z\)};
  \node[box,right=of cross] (r2) {\(\Rhyper\)};
  \node[boundary,right=16mm of r1,yshift=-15mm] (theta)
    {\(\Theta\) HMT canónico\\\scriptsize medida, carry, fase};
  \draw[->,thick,hmtblue] (tower)--(uhf);
  \draw[->,thick,hmtblue] (uhf)--(r1);
  \draw[->,thick,hmtblue] (clock)--(cross);
  \draw[->,thick,hmtblue] (cross)--(r2);
  \draw[<->,very thick,hmtgreen] (r1)--node[right]{\scriptsize isomorfía abstracta}(r2);
  \draw[dashed,->,hmtred] (r1)--(theta);
  \draw[dashed,->,hmtred] (r2)--(theta);
\end{tikzpicture}
}
\caption{Dos rutas abstractas hacia \(\Rhyper\). La conjugación causal HMT
es un problema adicional.}
\label{p12-83-c21-s3:fig:two-routes-r}
\end{figure}

La coincidencia de cierres débiles no identifica los antecedentes
\(C^*\). La rama UHF y el producto cruzado del odómetro pueden poseer
invariantes \(K\)-teóricos distintos. Una identificación HMT canónica exige
un mapa \(\Theta\) con
\[
  \Theta_*\mu_{\APP}=\mu_H,
  \qquad
  \Theta S_{\rm carry}=T_u\Theta,
\]
y preservación de diagonales, fase, orientación y proyectores. Ese
entrelazador permanece como frontera.

\subsection{Teorema ergódico medio de von Neumann para el odómetro de Haar en \(K_\infty\)}

Sea \(U f=f\circ T_u^{-1}\) en \(L^2(K_\infty,\mu_H)\). El teorema
ergódico medio de von Neumann \citep{vonneumann1932ergodic} da
\[
 \frac1N\sum_{j=0}^{N-1}U^jf
 \longrightarrow
 P_{\mathrm{Fix}(U)}f.
\]
La ergodicidad identifica el límite con
\[
  \left(\int f\,d\mu_H\right)\mathbf1.
\]
El reloj posee espectro puro puntual y entropía cero, pero sus promedios
temporales convergen al promedio de Haar. No se necesita mezcla.

\subsection{Retículo ortomodular de proyecciones y lógica de los subespacios cerrados}

El retículo \(\Proj(\Rhyper)\), con
\[
  P^\perp=I-P,
  \qquad P\wedge Q,
  \qquad P\vee Q=(P^\perp\wedge Q^\perp)^\perp,
\]
es completo y ortomodular. No es distributivo. Los proyectores
no conmutativos de APP se insertan en \(\Rhyper\) y proporcionan un testigo
finito de esa no distributividad \citep{birkhoffvonneumann1936}.

Cada subálgebra abeliana selecciona un contexto booleano. La lógica clásica
no desaparece: ocupa los contextos conmutativos. Lo que falla globalmente es
la posibilidad de reunir todos los observables incompatibles en una sola
álgebra booleana.

\subsection{Matrices de densidad y entropía}

En \(M_{9^d}\), un estado normal se escribe
\[
  \varphi_\rho(a)=\Tr(\rho a),
  \qquad \rho\ge0,\quad\Tr\rho=1.
\]
En \(\Rhyper\), una densidad normal es
\[
  h\in L^1(\Rhyper,\tau)_+,
  \qquad \tau(h)=1,
\]
y \(\varphi_h(a)=\tau(ha)\). Su entropía relativa a la traza es
\[
  S_\tau(h)=-\tau(h\log h).
\]

Si \(h=9^d\rho\), entonces
\[
  S_\tau(h)=S_{\rm vN}(\rho)-d\log9.
\]
La extensión canónica
\[
  \rho_{d+1}=\rho_d\otimes\frac{I_9}{9}
\]
añade \(\log9\) a la entropía matricial, pero deja invariante
\(S_\tau(h)\). La primera cantidad registra una fibra uniforme adicional;
la segunda usa la normalización holográfica de la torre.

\begin{controlbox}
No se usarán indistintamente \(\Tr\) y \(\tau\). Tampoco se interpretará
\(\tau(P)\) como cardinal de puntos, ni se llamará puro a todo vector de una
representación. Un factor difuso \(II_1\) no posee proyecciones mínimas.
\end{controlbox}

\subsection{Tipos operatoriales y frontera modular}

La clasificación construida hasta aquí es:
\[
\begin{array}{c|c}
\text{objeto}&\text{tipo}\\
\hline
M_{9^d}&I_{9^d}\\
\text{rama marcada}&I_\infty\\
\text{rama UHF tracial}&II_1\\
\text{reloj p.m.p.}&II_1.
\end{array}
\]
Un tipo \(III\) exigiría una medida cuasi-invariante con cociclo de
Radon--Nikodym no trivial o un estado KMS cuya representación GNS no fuese
semifinita. La mera presencia de una palabra «KMS» no selecciona el tipo.

\enlargethispage{6\baselineskip}
\begin{workbox}
La completación del grupoide total de memoria del \TPK\ requiere topología
o estructura boreliana, medida, cierre bajo inversión, amenabilidad,
ergodicidad y control de isotropía. El subgrupoide del reloj satisface esas
hipótesis; el grupoide total no se promociona por analogía.
\end{workbox}
